\chapter{Lines, Circles and Angles}\label{ch:g6:lines}

Geometry starts with a ruler, a set square and a compass. This chapter
fixes the vocabulary --- points, segments, rays, lines, circles ---
introduces the two special positions of lines (parallel and
perpendicular), and teaches how to measure and draw angles.

\section{Points, segments, lines}

\begin{definition}[Basic objects]\label{def:g6:lines:objects}
Through two distinct points $A$ and $B$ pass:
\begin{itemize}
  \item the \emph{segment}\index{segment} $[AB]$: the part of the line
        between $A$ and $B$ (it has a length, written $AB$);
  \item the \emph{ray} $[AB)$: starts at $A$, goes through $B$ and
        continues forever;
  \item the \emph{line}\index{line} $(AB)$: extends forever on both
        sides. Two points determine exactly one line.
\end{itemize}
Points on the same line are called \emph{aligned}.
\end{definition}

\begin{omfigure}
\begin{tikzpicture}[scale=1.0]
  \fill (0,1.8) circle (1.5pt) node[above, font=\small] {$A$};
  \fill (2.5,1.8) circle (1.5pt) node[above, font=\small] {$B$};
  \draw[thick, omDef] (0,1.8) -- (2.5,1.8);
  \node[right, font=\small] at (2.9,1.8) {segment $[AB]$};
  \fill (0,0.9) circle (1.5pt) node[above, font=\small] {$A$};
  \fill (2.5,0.9) circle (1.5pt) node[above, font=\small] {$B$};
  \draw[thick, omThm] (0,0.9) -- (4.1,0.9);
  \node[right, font=\small] at (4.3,0.9) {ray $[AB)$};
  \fill (0,0) circle (1.5pt) node[above, font=\small] {$A$};
  \fill (2.5,0) circle (1.5pt) node[above, font=\small] {$B$};
  \draw[thick, omMeth] (-1.4,0) -- (4.1,0);
  \node[right, font=\small] at (4.3,0) {line $(AB)$};
\end{tikzpicture}

{\small Same two points, three different objects. The brackets say what
stops and what continues: $[$ stops, $($ continues.}
\end{omfigure}

\section{Parallel and perpendicular lines}

\begin{definition}[Perpendicular, parallel]\label{def:g6:lines:perp}
Two lines are \emph{perpendicular}\index{perpendicular lines} when they
cross at a right angle; we write $d_1 \perp d_2$ and mark the right
angle with a small square. Two lines are
\emph{parallel}\index{parallel lines} when they never meet, however far
they are extended; we write $d_1 \parallel d_2$.
\end{definition}

\begin{omfigure}
\begin{tikzpicture}[scale=1.0]
  \draw[thick] (-1.8,0) -- (1.8,0) node[right, font=\small] {$d_1$};
  \draw[thick, omDef] (0.6,-1.2) -- (0.6,1.4)
    node[above, font=\small] {$d_2$};
  \draw[thin] (0.6,0) ++(-0.25,0) -- ++(0,0.25) -- ++(0.25,0);
  \node[below, font=\small] at (0,-1.5) {$d_1 \perp d_2$};
  \begin{scope}[xshift=6.5cm]
  \draw[thick] (-1.8,-0.5) -- (1.8,0.3) node[right, font=\small] {$d_3$};
  \draw[thick, omThm] (-1.8,-1.3) -- (1.8,-0.5)
    node[right, font=\small] {$d_4$};
  \node[below, font=\small] at (0,-1.5) {$d_3 \parallel d_4$};
  \end{scope}
\end{tikzpicture}

{\small A right-angle crossing (marked with the little square) and two
parallels: same direction, no crossing point.}
\end{omfigure}

\begin{proposition}[Two useful facts]\label{prop:g6:lines:facts}
\begin{enumerate}
  \item If two lines are both perpendicular to a third line, they are
        parallel to each other.
  \item If two lines are parallel, every line perpendicular to one is
        perpendicular to the other.
\end{enumerate}
\end{proposition}
\admitted

\begin{method}[Drawing with the set square]\label{met:g6:lines:setsquare}
To draw the perpendicular to a line $d$ through a point $P$:
\begin{enumerate}
  \item place one edge of the right angle of the set square along $d$;
  \item slide the set square along $d$ until its other edge reaches
        $P$;
  \item draw the line along that edge, and mark the right angle.
\end{enumerate}
For a parallel through $P$: draw a perpendicular to $d$, then the
perpendicular to \emph{that} line through $P$
(\cref{prop:g6:lines:facts}).
\end{method}

\section{Circles}

\begin{definition}[Circle]\label{def:g6:lines:circle}
The \emph{circle}\index{circle} of center $O$ and radius $r$ is the set
of all points at distance exactly $r$ from $O$. A segment from the
center to the circle is a \emph{radius}; a segment joining two points
of the circle through the center is a \emph{diameter} --- its length is
$2r$; a segment joining two points of the circle is a \emph{chord}.
\end{definition}

\begin{omfigure}
\begin{tikzpicture}[scale=1.15]
  \draw[thick] (0,0) circle (1.5);
  \fill (0,0) circle (1.5pt) node[below left, font=\small] {$O$};
  \fill (1.5,0) circle (1.5pt) node[right, font=\small] {$M$};
  \draw[omDef, thick] (0,0) -- (1.5,0)
    node[midway, above, font=\small] {radius};
  \fill (-1.06,1.06) circle (1.5pt) node[above left, font=\small] {$A$};
  \fill (1.06,-1.06) circle (1.5pt) node[below right, font=\small] {$B$};
  \draw[omThm, thick] (-1.06,1.06) -- (1.06,-1.06)
    node[pos=0.25, right=2pt, font=\small] {diameter};
  \fill (0.52,1.41) circle (1.5pt) node[above, font=\small] {$C$};
  \fill (1.41,0.52) circle (1.5pt) node[right, font=\small] {$D$};
  \draw[omMeth, thick] (0.52,1.41) -- (1.41,0.52)
    node[midway, above right=-2pt, font=\small] {chord};
\end{tikzpicture}

{\small A circle with a radius $[OM]$, a diameter $[AB]$ (a chord
through the center, twice as long as the radius) and a chord $[CD]$.}
\end{omfigure}

\begin{example}\label{ex:g6:lines:circle}
``Draw the circle of center $O$ passing through $A$'': open the compass
from $O$ to $A$ --- the radius is the distance $OA$ --- and turn. Every
point of this circle is at the same distance from $O$ as $A$.
\end{example}

\section{Angles}

\begin{definition}[Angle]\label{def:g6:lines:angle}
Two rays $[AB)$ and $[AC)$ with the same starting point $A$ form the
\emph{angle}\index{angle} $\widehat{BAC}$; the point $A$ is its
\emph{vertex} (always the middle letter!). Angles are measured in
\emph{degrees} ($^\circ$), from $0^\circ$ to $360^\circ$ for a full
turn. An angle is:
\begin{itemize}
  \item \emph{right} if it measures $90^\circ$;
  \item \emph{acute} if it measures less than $90^\circ$;
  \item \emph{obtuse} if it measures between $90^\circ$ and
        $180^\circ$;
  \item \emph{straight} if it measures $180^\circ$ (the two rays form
        a line).
\end{itemize}
\end{definition}

\begin{omfigure}
\begin{tikzpicture}[scale=0.9]
  \draw[thick] (0,0) -- (1.8,0);
  \draw[thick] (0,0) -- (40:1.8);
  \draw[omDef, ->] (0.7,0) arc (0:40:0.7);
  \node[omDef, font=\small] at (20:1.0) {$40^\circ$};
  \node[below, font=\small] at (0.9,-0.3) {acute};
  \begin{scope}[xshift=3.6cm]
  \draw[thick] (0,0) -- (1.8,0);
  \draw[thick] (0,0) -- (0,1.8);
  \draw[thin] (0.3,0) -- (0.3,0.3) -- (0,0.3);
  \node[below, font=\small] at (0.9,-0.3) {right};
  \end{scope}
  \begin{scope}[xshift=7.2cm]
  \draw[thick] (0,0) -- (1.8,0);
  \draw[thick] (0,0) -- (135:1.8);
  \draw[omThm, ->] (0.6,0) arc (0:135:0.6);
  \node[omThm, font=\small] at (65:1.0) {$135^\circ$};
  \node[below, font=\small] at (0.9,-0.3) {obtuse};
  \end{scope}
  \begin{scope}[xshift=11.6cm]
  \draw[thick] (-1.5,0) -- (1.5,0);
  \fill (0,0) circle (1.2pt);
  \draw[omMeth, ->] (0.5,0) arc (0:180:0.5);
  \node[below, font=\small] at (0,-0.3) {straight};
  \end{scope}
\end{tikzpicture}

{\small The four families of angles. The right angle ($90^\circ$) is
the reference: acute means smaller, obtuse means larger.}
\end{omfigure}

\begin{method}[Measuring an angle with a protractor]\label{met:g6:lines:protractor}
\begin{enumerate}
  \item Place the center of the protractor exactly on the vertex of the
        angle;
  \item align its $0^\circ$ line with one side of the angle;
  \item read the graduation crossed by the other side --- using the
        scale that starts at $0$ on the aligned side;
  \item sanity-check with the eye: an acute angle must read less than
        $90^\circ$, an obtuse one more.
\end{enumerate}
\end{method}

\begin{example}\label{ex:g6:lines:estimate}
Before measuring, estimate! An angle slightly more open than the corner
of a sheet of paper is a little over $90^\circ$; half a right angle is
$45^\circ$; a third of a right angle is $30^\circ$. If your protractor
says $150^\circ$ for an angle that looks acute, you read the wrong
scale: the correct measure is $180^\circ - 150^\circ = 30^\circ$.
\end{example}

\section{Exercises}

\begin{exercise}[$\star$]\label{exo:g6:lines:1}
Draw three points $A$, $B$, $C$ not aligned. Draw in different colors:
the segment $[AB]$, the ray $[CA)$, the line $(BC)$.
\end{exercise}

\begin{exercise}[$\star$]\label{exo:g6:lines:2}
True or false? ``$[AB]$ and $[BA]$ are the same segment.''
``$[AB)$ and $[BA)$ are the same ray.'' ``$(AB)$ and $(BA)$ are the
same line.'' Explain each answer.
\end{exercise}

\begin{exercise}[$\star$]\label{exo:g6:lines:3}
Draw a line $d$ and a point $P$ not on $d$. Construct with the set
square: the perpendicular to $d$ through $P$, then the parallel to $d$
through $P$. Describe your steps.
\end{exercise}

\begin{exercise}[$\star$]\label{exo:g6:lines:4}
Lines $a$ and $b$ are both perpendicular to a line $c$, and a fourth
line $e$ is perpendicular to $a$. What can you say about $a$ and $b$?
About $e$ and $c$? Justify with \cref{prop:g6:lines:facts}.
\end{exercise}

\begin{exercise}[$\star$]\label{exo:g6:lines:5}
Draw a circle of center $O$ with radius $3$ cm. Place a point $M$ on
the circle, a point $N$ inside, a point $P$ outside. What can you say
about the distances $OM$, $ON$, $OP$ compared with $3$ cm?
\end{exercise}

\begin{exercise}[$\star$]\label{exo:g6:lines:6}
A circle has diameter $9$ cm. What is its radius? Another has radius
$2.6$ cm: what is its diameter?
\end{exercise}

\begin{exercise}[$\star$]\label{exo:g6:lines:7}
Name the marked angle in three letters, then classify it (acute, right,
obtuse, straight): an angle of $72^\circ$ at vertex $R$ between rays
towards $S$ and $T$; an angle of $148^\circ$ at vertex $B$ between rays
towards $A$ and $C$; an angle of $90^\circ$ at vertex $O$ between rays
towards $M$ and $N$.
\end{exercise}

\begin{exercise}[$\star$]\label{exo:g6:lines:8}
Estimate, then measure with a protractor, the three angles of a
triangle you draw yourself. Add the three measures: what do you find?
(Keep your answer for \cref{ch:g7:triangles}.)
\end{exercise}

\begin{exercise}[$\star$]\label{exo:g6:lines:9}
Draw an angle $\widehat{xOy}$ of $65^\circ$ with a protractor, then an
angle of $130^\circ$. How could you get the second one from the first
without the protractor?
\end{exercise}

\begin{exercise}[$\star\star$]\label{exo:g6:lines:10}
Two villages $A$ and $B$ are drawn on a map. Where are the points that
are at $3$ cm from $A$ \emph{and} at $2$ cm from $B$? Draw a picture
showing how many such points there can be ($0$, $1$ or $2$ depending on
the distance $AB$).
\end{exercise}

\begin{exercise}[$\star\star$]\label{exo:g6:lines:11}
A clock shows 3 o'clock: what is the angle between the two hands? Same
question at 5 o'clock, and at 6 o'clock. (A full turn is $360^\circ$
for $12$ hours.)
\end{exercise}

\begin{exercise}[$\star\star\star$]\label{exo:g6:lines:12}
Draw a segment $[AB]$ of $6$ cm. Construct the point $C$ such that
$AC = BC = 6$ cm, using only the compass, and measure the angle
$\widehat{CAB}$. What triangle did you build, and what do you
conjecture about its angles?
\end{exercise}
